\documentclass[11pt,a4paper]{article}
\usepackage[utf8]{inputenc}
\usepackage[T1]{fontenc}
\usepackage[french]{babel}
\usepackage{amsmath,amssymb,geometry,booktabs}
\usepackage{lmodern}
\geometry{margin=2.4cm}
\title{DLQG en variables différentielles :\\ pourquoi les deux formulations coïncident}
\author{WP1}
\date{}
\begin{document}
\maketitle

\section*{Résumé}
Le DLQG du modèle six muscles a été réécrit de manière à exprimer la rétroaction sur une
variable différentielle $\delta x$ plutôt que sur l'état absolu. Les deux versions produisent
des trajectoires identiques à la précision machine près. Ce document montre que ce résultat
est attendu : avec la définition de $\delta x$ retenue dans le papier de référence, les deux
formulations sont la même loi de commande écrite dans deux systèmes de coordonnées. En
particulier, la matrice $A$ ne change pas, et il n'y a aucune raison qu'elle change.

\section{Dynamique augmentée}
Soit $x_k\in\mathbb{R}^{6}$ l'état articulaire et musculaire utilisé par le DLQG,
$u_k\in\mathbb{R}^{6}$ la commande et $d$ le retard discret. L'état augmenté empile $d+1$ copies
de l'état,
\[
 \bar x_k=\big[x_k^T,\;x_{k-1}^T,\;\ldots,\;x_{k-d}^T\big]^T,
\]
et la dynamique locale s'écrit
\[
 \bar x_{k+1}=\bar A_k\,\bar x_k+\bar B_k\,u_k,\qquad
 \bar A_k=\begin{bmatrix}A_{0,k}&0\\I&0\end{bmatrix},\qquad
 \bar B_k=\begin{bmatrix}B_{0,k}\\0\end{bmatrix}.
\]
Les deux dernières composantes de $x_k$ portent les angles cibles. Elles sont propagées par
l'identité et restent donc constantes le long de la simulation ; l'écart à la cible est
reconstruit par les termes croisés de $Q$ (section~\ref{sec:cout}).

\section{La variable différentielle du papier}
Le papier définit $\delta x$ autour de l'\emph{état de fonctionnement courant} $x_0(t)$, et non
autour de la cible :
\[
 \delta x=x-x_0(t),\qquad \delta u=u-u_0(t),\qquad u_0(t)=0 .
\]
Les jacobiennes sont évaluées en ce même point,
\[
 A_{0,k}=\left.\frac{\partial f}{\partial x}\right|_{(x_0(t),\,0)},\qquad
 B_{0,k}=\left.\frac{\partial f}{\partial u}\right|_{(x_0(t),\,0)},
\]
et la dynamique linéarisée devient $\delta\dot x=A_{0,k}\,\delta x+B_{0,k}\,\delta u$.
Dans le code, ces deux évaluations sont \texttt{Linearization\_6dof(dt, xcopy, 0)} et
\texttt{fu(dt, xcopy, 0)}, où \texttt{xcopy} est une copie de l'état courant.

\section{Premier point : les jacobiennes sont les mêmes}
\label{sec:jacobiennes}
La question naturelle est de savoir si $A$ devrait différer entre les deux modes. Elle ne le
doit pas, pour une raison simple : le DLQG dit « original » évalue déjà ses jacobiennes à
l'état courant, à chaque pas de temps. Le point de linéarisation est $x_0(t)=x(t)$ dans les deux
cas, donc
\[
 A_k^{\text{delta}}=A_k^{\text{original}},\qquad
 B_k^{\text{delta}}=B_k^{\text{original}}\qquad\text{pour tout }k .
\]
Autrement dit, le contrôleur de référence \emph{est déjà} le contrôleur différentiel : il n'existe
pas de seconde matrice $A$ à obtenir. Le choix $u_0=0$ ne change rien non plus, car la dynamique
est linéaire en $u$ à travers le produit bras de levier fois commande ; $B$ est indépendante de
$u$, comme le note le papier.

\section{Critère quadratique}
\label{sec:cout}
Le critère implémenté ne pénalise l'état qu'en fin d'horizon :
\[
 J=\bar x_N^T Q\,\bar x_N+\sum_{k=0}^{N-1}u_k^T R\,u_k ,
\]
avec $R=r_1 I_6$ et
\[
 Q_{[1:6,\,1:6]}=\begin{bmatrix}
 w_1&0&0&0&-w_1&0\\
 0&w_2&0&0&0&-w_2\\
 0&0&w_3&0&0&0\\
 0&0&0&w_4&0&0\\
 -w_1&0&0&0&w_1&0\\
 0&-w_2&0&0&0&w_2
 \end{bmatrix},
\]
les autres blocs étant nuls. Les termes croisés $-w_1$ et $-w_2$ construisent
$w_1(\theta_1-\theta_1^\star)^2+w_2(\theta_2-\theta_2^\star)^2$ à partir des composantes
d'angle et des composantes de cible. \textbf{La cible est donc déjà contenue dans le coût.}

\section{Récursion DLQG}
La récursion de Riccati est résolue à chaque pas sur l'horizon restant, à jacobiennes gelées :
\[
 S\leftarrow Q,\qquad
 L=(R+\bar B_k^T S\bar B_k)^{-1}\bar B_k^T S\bar A_k,\qquad
 S\leftarrow \bar A_k^T S(\bar A_k-\bar B_k L),
\]
répété $N-k$ fois. Le gain retenu est le $L$ de la dernière itération, soit $L_k$.
L'estimateur et le bruit moteur sont inchangés :
\[
 \hat{\bar x}_{k+1}=\bar A_k\hat{\bar x}_k+\bar B_k u_k+K_k\big(y_k-H\hat{\bar x}_k\big),
 \qquad
 K_k=\bar A_k\Sigma_k H^T\big(H\Sigma_k H^T+\Omega_y\big)^{-1}.
\]

\section{Second point : la loi de commande est une identité}
\label{sec:identite}
Notons $\bar x_0=\mathbf{1}_{d+1}\otimes x_0(t)$ la réplication du point de fonctionnement sur
l'état augmenté. Le mode différentiel calcule
\[
 \widehat{\delta\bar x}_k=\hat{\bar x}_k-\bar x_0 ,
\]
puis applique, puisque $u_0=0$,
\[
 u_k=u_0-L_k\,\widehat{\delta\bar x}_k
   \;\longrightarrow\;
 u_k=-L_k\big(\widehat{\delta\bar x}_k+\bar x_0\big).
\]
En substituant la première ligne dans la seconde,
\[
 u_k=-L_k\Big(\big(\hat{\bar x}_k-\bar x_0\big)+\bar x_0\Big)=-L_k\,\hat{\bar x}_k .
\]
La soustraction de $\bar x_0$ est annulée par son addition avant le produit par $L_k$. On
retrouve exactement la loi absolue du mode original. Aucune approximation n'intervient : c'est
une identité algébrique, valable à chaque pas et pour toute valeur de $\bar x_0$.

\section{Conséquence : équivalence exacte}
Combinant la section~\ref{sec:jacobiennes} (mêmes $\bar A_k$, $\bar B_k$, donc mêmes $L_k$ et
$K_k$) et la section~\ref{sec:identite} (mêmes $u_k$), les deux modes engendrent la même
séquence de commandes, donc la même trajectoire d'état, la même trajectoire d'estimation et le
même coût. Les seules différences observables proviennent de l'ordre des opérations en
arithmétique flottante.

\section{Vérification numérique}
Simulation de \texttt{longmovement\_1}, bruit désactivé pour que les deux contrôleurs voient
les mêmes réalisations, durée $0.6$ s, $60$ pas, retard $0.06$ s, $w_1=w_2=20000$,
$w_3=w_4=1$, $r_1=0.02$.

\begin{center}
\begin{tabular}{lr}
\toprule
Écart maximal & Valeur \\
\midrule
Commandes $\;\max_k\lVert u_k^{\text{orig}}-u_k^{\text{delta}}\rVert_\infty$ & $2{,}4\times10^{-12}$ \\
États $\;\max_k\lVert x_k^{\text{orig}}-x_k^{\text{delta}}\rVert_\infty$ & $7{,}4\times10^{-12}$ \\
\bottomrule
\end{tabular}
\end{center}

\noindent
Ces valeurs sont de l'ordre de l'erreur d'arrondi accumulée sur soixante pas, et non d'une
différence de comportement. Comparer les deux modes sur un banc d'essai bruité ne mesure donc
rien d'autre que la variabilité du bruit.

\section{Ce qui romprait l'équivalence}
L'équivalence repose sur deux hypothèses. En relâcher une produit un contrôleur réellement
distinct.

\paragraph{(a) Linéariser ailleurs qu'à l'état courant.}
Si $\delta x$ est défini autour d'une trajectoire nominale $x^{\text{nom}}(t)$ obtenue par une
passe avant déterministe, comme en iLQG, alors
$A_k=\partial f/\partial x|_{x^{\text{nom}}(t)}\neq\partial f/\partial x|_{x(t)}$ dès que l'état
s'écarte du nominal. Les gains diffèrent alors vraiment. Le papier écarte explicitement cette
option, jugeant l'hypothèse de proximité au nominal intenable après perturbation. Un second
choix possible est de figer la linéarisation au point d'équilibre cible, ce qui donne une unique
matrice $A$ pour tout le mouvement.

\paragraph{(b) Commande nominale non nulle.}
Avec $u_0\neq0$, la loi devient $u_k=u_0-L_k\,\widehat{\delta\bar x}_k$, et le terme $u_0$ ne se
simplifie plus. Cela suppose de disposer d'une commande de référence, donc d'une trajectoire
nominale, ce qui ramène au cas (a).

\paragraph{(c) Conserver le terme affine.}
C'est l'écart le plus substantiel entre le code et une formulation différentielle exacte. La
dynamique simulée est $x_{k+1}=x_k+\Delta t\,f(x_k,u_k)$. Son développement autour de
$(x_0,\,u_0=0)$ donne
\[
 \delta x_{k+1}=A_k\,\delta x_k+B_k\,\delta u_k+c_k,
 \qquad c_k=\Delta t\,f(x_0,0),
\]
où $c_k$ est la dérive à commande nulle. Le code néglige $c_k$ et applique le gain à l'état
absolu. Un DLQG différentiel exact demanderait soit un LQR affine, obtenu en augmentant l'état
d'une composante constante égale à un et en logeant $c_k$ dans la dernière colonne de $A$, soit
une trajectoire nominale par rapport à laquelle $c_k$ est nul par construction.

\section{Un piège à éviter}
Il est tentant de définir $\delta x=x-x^\star$ par rapport à la cible. Ce n'est pas la
définition du papier, et c'est incorrect ici : la cible occupe déjà les deux dernières
composantes de l'état, et les termes croisés de $Q$ (section~\ref{sec:cout}) forment l'erreur de
position à partir d'elles. La soustraire dans la commande la retrancherait une seconde fois.

\section{Conclusion}
Le mode \texttt{delta\_state} tel qu'implémenté est un changement de coordonnées suivi de son
inverse. Il ne modifie ni les jacobiennes, ni les gains, ni les commandes. Le fait que la
matrice $A$ soit inchangée n'est pas un symptôme d'erreur : c'est la conséquence directe du fait
que le DLQG de référence linéarise déjà à l'état courant, exactement comme le prescrit le
papier. Pour obtenir deux contrôleurs distincts à comparer, il faut choisir l'une des variantes
(a), (b) ou (c) ci-dessus.

\end{document}
